\section{Data: filter cascade, descriptor classes and integrity exclusions}
\label{si:data}

The dataset of this work is the set of neutral oxygen vacancies, $V_\mathrm{O}^{\times}$, of the release of
Kumagai et al.\ \citep{Kumagai2021insights} that survive a fixed cascade of filters (Table~\ref{tab:cascade}; the
cascade is drawn as Figure~\ref{fig:cascade} in Section~\ref{si:splits}). The release holds
\numFilterCascadeStepsReleaseSites{} completed neutral calculations in \numFilterCascadeStepsReleaseHosts{} host
oxides. Kumagai et al.\ built from them a neutral machine-learning subset of \numFilterCascadeStepsMlSubsetSites{}
vacancies in \numFilterCascadeStepsMlSubsetHosts{} hosts, from which they had already removed the
perturbed-host-state entries (\numScreeningAbsenceReleaseNeutralCompletedPhs{} of the completed neutral entries),
the entries whose band edges could not be determined and the dynamically unstable hosts. A perturbed-host-state
entry is a vacancy whose two donor electrons do not stay in a localised in-gap state but delocalise into the host
conduction band, so the subset, and every result of this work, describes deep, localised neutral vacancies in
ordered oxides classified as non-magnetic in the release. Our three further filters, D1 to D3, remove what the data audit
(\texttt{docs/data\_audit.md}) found inconsistent in that subset; they leave \numFilterCascadeStepsFinalSites{}
oxygen sites in \numFilterCascadeStepsFinalHosts{} hosts, one host per formula. The site counts after each
step are read from \texttt{results/filter\_cascade.json}.

\begin{table}[htbp]
\caption{Filter cascade of the \numFilterCascadeStepsFinalHosts{}-host dataset, with the oxygen sites, hosts and formulas left after each step
(\texttt{results/filter\_cascade.json}; the same table as \arttabcascade{} of the article).}
\label{tab:cascade}
\centering\small
\begin{adjustbox}{max width=\linewidth}
\begin{tabular}{llrrr}
\toprule
Step & Criterion & Sites & Hosts & Formulas \\
\midrule
Release & Completed neutral entries in the release & 2103 & 935 & 935 \\
ML subset & Kumagai neutral ML subset (\texttt{charge0.csv}) & 1745 & 824 & 824 \\
D1 & Vacancy split, unknown type or not same config from init & 1730 & 819 & 819 \\
D2 & Host Sr$_{2}$Zr$_{7}$O$_{16}$ (ML site labels do not match release targets) & 1726 & 818 & 818 \\
D3 & Ba$_{4}$Nb$_{2}$WO$_{12}$, Sr$_{2}$SnO$_{4}$: vacancy site is O with matching first cation shell & 1726 & 818 & 818 \\
\midrule
Final & Final universe & 1726 & 818 & 818 \\
\bottomrule
\end{tabular}

\end{adjustbox}
\end{table}

The three exclusions concern the integrity of the labels, not the chemistry of the hosts. D1 removes the
\numFilterCascadeStepsDOneNRemoved{} entries of the subset that the release itself flags as a split vacancy, an
unknown defect type, or a relaxed configuration that is not the configuration of the initial vacancy. Each of
them also carries the last flag, so the stored energy of these entries may not belong to a single vacancy on the
labelled site, which is the structure that a model reads. The subset had kept these entries although the release
flags them. D2 removes the
\numFilterCascadeStepsDTwoNRemoved{} sites of Sr$_2$Zr$_7$O$_{16}$, whose site labels in the machine-learning
tables do not match the targets of the release; the release gives no criterion that explains the mismatch, so the
host is excluded rather than repaired. D3 concerns the two hosts whose cell metadata are inconsistent with their
supercells, Ba$_4$Nb$_2$WO$_{12}$ and Sr$_2$SnO$_4$ (the Sr$_2$SnO$_4$ supercell of
\numStructureAuditAnomalousHostsSrTwoSnOFourNAtomsSc{} atoms has \numStructureAuditAnomalousHostsSrTwoSnOFourVolumeRatio{}
of the expected volume, because it was built from the primitive cell). An entry of these hosts is kept only if
its labelled vacancy site maps to an oxygen atom of the supercell whose first cation shell matches that of the
unit-cell site within \numFilterCascadeDThreeSrTwoSnOFourSitesOOneToleranceA{}~\AA{}; every entry passed, and D3
removed \numFilterCascadeStepsDThreeNRemoved{} sites. Every site of the dataset lies in the published split of
Kiyohara et al.\ \citep{Kiyohara2025machinelearning} (\numFilterCascadeKiyoharaSplitSitesAbsent{} sites absent),
which is used for the comparison with the published accuracy (Section~\ref{si:splits}).

The descriptors of the release fall into three classes (Table~\ref{tab:taxonomy}). The host-electronic class
holds properties of the pristine host from its DFT electronic structure: the band gap of the PBEsol(+$U$)
calculation, the electronic and ionic dielectric constants and the orbital character of the band edges, together
with the formation energy of the host oxide, which measures how strongly oxygen is bound. The site-electronic class
describes the oxygen atom to be removed in the pristine host: its Bader charge and volume, its Born effective
charge, the same quantities averaged over and extremal among its Voronoi neighbours, and the centre of its
O-2$p$ band below the valence-band maximum. The structural and compositional class holds the fractions of the
solid angle of the oxygen Voronoi cell that face each element and the electronegativities of the neighbours. D
receives the two electronic classes, \numRfElectronicNFeatures{} descriptors; RF-Kumagai uses all
\numRfKumagaiNFeatures{} descriptors, RF-electronic the \numRfElectronicNFeatures{} electronic ones and
RF-structural the \numRfStructuralNFeatures{} structural and compositional ones. Every descriptor is computed on
the pristine host, so none of them contains the relaxed defect geometry or the defect energy. The host formation
energy is named \texttt{formation\_energy} in the release and is not the vacancy formation energy that the models
predict.

{\footnotesize\tighttabs
% Generated by paper/si/tools/make_si_tables.py from docs/descriptor_taxonomy.csv. Do not edit.
\begin{longtable}{@{}>{\raggedright}p{0.25\linewidth}>{\raggedright}p{0.15\linewidth}>{\raggedright}p{0.135\linewidth}>{\raggedright\arraybackslash}p{0.385\linewidth}@{}}
\caption{The descriptors of the release of Kumagai et al.\ \citep{Kumagai2021insights} and their classes, transcribed from \texttt{docs/descriptor\_taxonomy.csv} (names as in the release; definition abridged to its first sentence). The host-electronic and site-electronic DFT descriptors are those given to D; RF-Kumagai uses all of them, RF-electronic the two electronic classes and RF-structural the structural and compositional class. Every descriptor is a property of the pristine host; \texttt{formation\_energy} is the formation energy of the host oxide, not the vacancy target.}\label{tab:taxonomy}\\
\toprule
Name & Class & Units & Definition \\
\midrule
\endfirsthead
\toprule
Name & Class & Units & Definition \\
\midrule
\endhead
\bottomrule
\endlastfoot
\texttt{ave\_ele\_diele} & host-electronic & dimensionless & Spherically averaged electronic (ion-clamped) dielectric constant of the pristine host, from DFPT \\
\texttt{ave\_ion\_diele} & host-electronic & dimensionless & Spherically averaged ionic dielectric constant of the pristine host, from DFPT \\
\texttt{band\_gap} & host-electronic & eV & Band gap of the pristine host from the PBEsol(+U) calculation; not the nsc-dd hybrid gap \\
\texttt{cbm\_d} & host-electronic & dimensionless & Fraction of d-orbital character in the pristine host's conduction-band minimum state (projection ratio) \\
\texttt{cbm\_f} & host-electronic & dimensionless & Fraction of f-orbital character in the pristine host's conduction-band minimum state (projection ratio) \\
\texttt{cbm\_p} & host-electronic & dimensionless & Fraction of p-orbital character in the pristine host's conduction-band minimum state (projection ratio) \\
\texttt{cbm\_s} & host-electronic & dimensionless & Fraction of s-orbital character in the pristine host's conduction-band minimum state (projection ratio) \\
\texttt{formation\_energy} & host-electronic & eV/atom & Host formation energy per atom retrieved from the Materials Project database (empirically corrected) \\
\texttt{vbm\_d} & host-electronic & dimensionless & Fraction of d-orbital character in the pristine host's valence-band maximum state (projection ratio) \\
\texttt{vbm\_f} & host-electronic & dimensionless & Fraction of f-orbital character in the pristine host's valence-band maximum state (projection ratio) \\
\texttt{vbm\_p} & host-electronic & dimensionless & Fraction of p-orbital character in the pristine host's valence-band maximum state (projection ratio) \\
\texttt{vbm\_s} & host-electronic & dimensionless & Fraction of s-orbital character in the pristine host's valence-band maximum state (projection ratio) \\
\addlinespace
\texttt{ave\_bec} & site-electronic & e & Spherically averaged Born effective charge (trace/3) of the O atom to be removed, in the pristine host (DFPT) \\
\texttt{bader\_charge} & site-electronic & e & Bader charge of the O atom to be removed, from the pristine host unit-cell all-electron density; net ionic charge, negative for the anion \\
\texttt{bader\_volume} & site-electronic & \AA$^{3}$ & Bader (zero-flux) volume of the O atom to be removed, from the pristine host all-electron density \\
\texttt{nn\_ave\_ave\_bec} & site-electronic & e & Solid-angle-weighted mean of the spherically averaged Born effective charge over the Voronoi neighbours \\
\texttt{nn\_ave\_bader\_charge} & site-electronic & e & Solid-angle-weighted mean of the Bader charge over the Voronoi neighbours of the O site \\
\texttt{nn\_max\_ave\_bec} & site-electronic & e & Maximum of the spherically averaged Born effective charge over the Voronoi neighbours \\
\texttt{nn\_max\_bader\_charge} & site-electronic & e & Maximum of the Bader charge over the Voronoi neighbours of the O site \\
\texttt{nn\_min\_ave\_bec} & site-electronic & e & Minimum of the spherically averaged Born effective charge over the Voronoi neighbours \\
\texttt{nn\_min\_bader\_charge} & site-electronic & e & Minimum of the Bader charge over the Voronoi neighbours of the O site \\
\texttt{o2p\_center\_from\_vbm} & site-electronic & eV & Centroid of the O-2p projected density of states of the O site, measured as its distance below the VBM (Deml et al \\
\addlinespace
\texttt{nn\_ave\_eleneg} & structural & Pauling units & Solid-angle-weighted mean of the Pauling electronegativity over the Voronoi neighbours (cations and O) \\
\texttt{nn\_max\_eleneg} & structural & Pauling units & Maximum of the Pauling electronegativity over the Voronoi neighbours (cations and O) \\
\texttt{nn\_min\_eleneg} & structural & Pauling units & Minimum of the Pauling electronegativity over the Voronoi neighbours (cations and O) \\
\texttt{w\_X} & structural & fraction & Fraction of the solid angle of the O-site Voronoi cell (pristine unit cell) that faces element X, for X in: Bi, Ag, O, Hg, Au, Ge, Al, Pb, Zn, As, Cu, Sb, Mo, Tl, B, Ba, Ca, Pd, Cd, Pt, Ce, Hf, In, Ga, La, Ta, Mg, Si, Sc, Sn, Ti, Y, Zr, Nb, Na, Sr, V, Li, Be, W, Cr, Cs, K, Re, Rb \\
\end{longtable}
}
